\subsection{Scene I --- The Companion and Grappling Bear}\label{subsec:scene_bear}

The first scene, \texttt{GrapplingBear}, places a single character --- a bear ---
under two distinct control regimes. \texttt{GameManager} holds a global grappling
flag, toggled by both hands' \texttt{Thumb\_Down}, and broadcasts it each frame to
every component that must respond. With the flag clear the bear is a
\emph{companion} that reacts to the user; with it set the bear becomes an
\emph{avatar} the user pilots. The receiver hides the hand meshes whenever the
flag is set, so the hands disappear the moment they take control of the bear.

As a companion the bear does two things. A \texttt{Thumb\_Up} from either hand
raises the bear's arm on the mirrored side and waves it, with a slight head tilt
and a dialogue bubble, handled by \texttt{BearInteraction}. Pointing at the ground
and confirming with the left hand's \texttt{Victory}, through the shared
\texttt{HandPointer}, sends the bear walking to that point. Throughout,
\texttt{CameraRingController} in its ring mode lets the user
orbit the scene by holding \texttt{Pointing\_Up} on one hand and \texttt{Closed\_Fist}
on the other, the camera tracing a circle about a fixed centre while preserving its
tilt.

\begin{figure}[H]
	\centering
	\captionsetup{justification=centering}
	\includegraphics[width=0.9\linewidth, frame]{img/scene_bear_companion.png}
	\caption{The companion bear waving back at the user
	         in response to a \texttt{Thumb\_Up}.}
	\label{fig:bear-companion}
\end{figure}

With the grappling flag set, the same bear is flown by the hands directly. Each
hand drives a virtual aim point offset from a centre fixed to the bear and scaled
from the hand's tracked position, and the bear's arms turn to follow these points.
A sphere cast along each arm's aim direction, thickened to forgive imprecise
pointing, reports whether a graspable surface lies ahead; closing that hand into a
\texttt{Closed\_Fist} on a valid target shoots a rope and attaches a spring joint
between the bear and the struck point. The joint, tuned with a stiff spring and
heavy damping and a rest length shorter than the initial throw, pulls the bear in
and lets it swing, with the bear's forward direction locked toward the anchor for
the duration. Releasing the fist detaches the rope, subject to the brief grace
window of~\autoref{subsec:interaction_layer}.

The two hands divide the labour of grappling and looking. When one hand holds a
rope the other is free, and \texttt{CameraRingController} switches to a third-person
view that the free hand steers: its tracked position turns the camera's yaw about
the bear, and the view is aimed at the point the free hand is itself pointing at.
One hand thus commits the bear to a swing while the other chooses where to look and
where to throw next.

A collision routine in the shared grapple core classifies each impact by comparing
its direction against the bear's own forward, right, and up axes, and, when an
optional haptic vest is connected, fires the motor bank on the corresponding side
of the torso. The vest is strictly an enhancement. The interaction itself needs
only the webcam pipeline that defines this thesis~(\autoref{subsec:gap}); the
haptic feedback, when present, adds a felt response to collisions in this one scene
and nothing more. Taken whole, the scene answers the games question
of~\autoref{subsubsec:sota_immersive}.
The two-handed grapple-and-swing it demonstrates is the kind of
hand-driven control shipped in immersive titles --- the same control
here, however, is driven from a single webcam.

\begin{figure}[H]
	\centering
	\captionsetup{justification=centering}
	\includegraphics[width=0.9\linewidth, frame]{img/scene_bear_grapple.png}
	\caption{Grappling mode: one hand holds a rope while the
	         free hand steers the third-person camera.}
	\label{fig:bear-grapple}
\end{figure}
